% Part: many-valued-logic
% Chapter: three-valued-logics
% Section: lukasiewicz

\documentclass[../../../include/open-logic-section]{subfiles}

\begin{document}

\olfileid{mvl}{thr}{luk}

\olsection{منطق~\L ukasiewicz}

يرجع أحد أوائل المقترحات المنشورة والمفصلة لمنطق متعدد القيم إلى
الفيلسوف البولندي يان~\L ukasiewicz سنة~1921. وكان~\L ukasiewicz
مدفوعًا بمسألة
أرسطو عن المعركة البحرية: يبدو أن الجملة ``ستقع معركة بحرية غدًا''
ليست، \emph{اليوم}، صادقة ولا كاذبة؛ إذ لم تتحدد قيمة صدقها بعد.
واقترح~\L ukasiewicz إدخال قيمة صدق ثالثة لمثل هذه الجمل
``المستقبلية العرضية''.
\begin{quote}
يمكنني أن أفترض، من دون تناقض، أن وجودي في وارسو في لحظة معينة من
السنة المقبلة، كظهر يوم~21 ديسمبر مثلًا، غير محدد الآن لا إيجابًا ولا
سلبًا. ولذلك فمن الممكن، من غير الضروري، أن أكون في وارسو في الوقت
المعين. وبناء على هذا الافتراض، لا تكون القضية ``سأكون في وارسو ظهر
يوم~21 ديسمبر من السنة المقبلة'' في الوقت الحاضر صادقة ولا كاذبة.
فلو كانت صادقة الآن، لكان وجودي المستقبلي في وارسو ضروريًا، وهذا يناقض
الافتراض. ولو كانت كاذبة الآن، من جهة أخرى، لكان وجودي المستقبلي في
وارسو مستحيلًا، وهذا يناقض الافتراض أيضًا. ومن ثم فالقضية موضع النظر
ليست الآن صادقة ولا كاذبة، ولا بد أن تكون لها قيمة ثالثة مختلفة عن
``$0$'' أو الكذب وعن ``$1$'' أو الصدق. ويمكننا أن نرمز إلى هذه القيمة
بـ``$\frac{1}{2}$''. وهي تمثل ``الممكن''، وتنضم إلى ``الصادق''
و``الكاذب'' قيمة ثالثة.
\end{quote}
سنستعمل~$\Undef$ لقيمة الصدق الثالثة عند~\L
ukasiewicz.\footnote{يستعمل~\L ukasiewicz كلمة ``ممكن'' هنا بمعنى
غير مألوف اليوم، هو الممكن غير الضروري.}

يسهل تحديد دوال الصدق للروابط~$\lnot$ و$\land$ و$\lor$ وفق هذا
التفسير: فنفي جملة مستقبلية عرضية جملة مستقبلية عرضية أيضًا، ولذلك
$\tf{\lnot}(\Undef) = \Undef$. وإذا كان أحد مقترني اقتران ما غير
محدد والآخر صادقًا، كان الاقتران غير محدد أيضًا؛ فقد يتبين أن الاقتران
صادق، وقد يتبين أنه كاذب، بحسب ما يؤول إليه المقترن المستقبلي العرضي.
ومن ثم
\[
    \tf{\land}(\True, \Undef) =
\tf{\land}(\Undef, \True) =
\Undef.
\]
أما إذا كان المقترن الآخر كاذبًا، فلا يمكن أن يتبين أن الاقتران صادق؛
لذلك
\[\tf{\land}(\False, \Undef) =
\tf{\land}(\Undef, \False) = \False.\]
أما القيم الأخرى (عندما تكون الوسائط قيم صدق متعينة، أي~$\True$ أو
$\False$) فهي كما في المنطق الكلاسيكي.

أما الشرط فحاله أعقد قليلًا. افترض أن~$q$ عبارة مستقبلية عرضية. إذا
كان~$p$ كاذبًا، كان~$p \lif q$ صادقًا بصرف النظر عما يؤول إليه~$q$؛
ولذلك ينبغي أن نضع~$\tf{\lif}(\False, \Undef) = \True$. وإذا كان~$p$
صادقًا، كان~$q \lif p$ صادقًا بصرف النظر عما يؤول إليه~$q$؛ ولذلك
$\tf{\lif}(\Undef, \True) = \True$. وإذا كان~$p$ صادقًا، فقد يتبين
أن~$p \lif q$ صادق أو كاذب؛ ولذلك
$\tf{\lif}(\True, \Undef) = \Undef$. وبالمثل، إذا كان~$p$ كاذبًا،
فقد يتبين أن~$q \lif p$ صادق أو كاذب؛ ولذلك
$\tf{\lif}(\Undef, \False) = \Undef$. وتبقى الحالة التي يكون فيها
$p$ و$q$ كلاهما مستقبليًا عرضيًا. وبحسب الدافع، ينبغي حقًا أن نسند
$\Undef$ في هذه الحالة. غير أن ذلك يجعل~$!A \lif !A$ \emph{غير}
تحصيل منطقي. ولم يجد~\L ukasiewicz مشقة في التخلي عن
$!A \lor \lnot !A$ و$\lnot(!A \land \lnot !A)$، لكنه تردد في التخلي
عن~$!A \lif !A$. ولذلك اشترط
$\tf{\lif}(\Undef, \Undef) = \True$.

\begin{defn}\ollabel{def:lukasiewicz}
يعرّف منطق~\L ukasiewicz الثلاثي القيم بالمصفوفة الآتية:
\begin{enumerate}
  \item لغة القضايا القياسية~$\Lang L_0$ بروابطها
  $\lnot$ و$\land$ و$\lor$ و$\lif$.
  \item مجموعة قيم الصدق~$V = \{\True, \Undef, \False\}$.
  \item $\True$ هي القيمة المميزة الوحيدة؛ أي~$V^+ = \{\True\}$.
  \item تعطي الجداول الآتية دوال الصدق:
  \begin{center}
    \begin{tabular}{c|c}
      $\tf{\lnot}$ & \\
      \hline
      $\True$ & $\False$ \\
      $\Undef$ & $\Undef$ \\
      $\False$ & $\True$
    \end{tabular}
    \quad
    \begin{tabular}{c|ccc}
      $\tf{\land}[\LogLuk[3]]$ & $\True$ & $\Undef$ & $\False$ \\
      \hline
      $\True$ & $\True$ & $\Undef$ & $\False$ \\
      $\Undef$ & $\Undef$ & $\Undef$ & $\False$\\
      $\False$ & $\False$ & $\False$ & $\False$
    \end{tabular}
    \\[2ex]
    \begin{tabular}{c|ccc}
      $\tf{\lor}[\LogLuk[3]]$ & $\True$ & $\Undef$ & $\False$ \\
      \hline
      $\True$ & $\True$ & $\True$ & $\True$ \\
      $\Undef$ & $\True$ & $\Undef$ & $\Undef$ \\
      $\False$ & $\True$ & $\Undef$ & $\False$
    \end{tabular}
    \quad
    \begin{tabular}{c|ccc}
      $\tf{\lif}[\LogLuk[3]]$ & $\True$ & $\Undef$ & $\False$ \\
      \hline
      $\True$ & $\True$ & $\Undef$ & $\False$ \\
      $\Undef$ & $\True$ & $\True$ & $\Undef$  \\
      $\False$ & $\True$ & $\True$ & $\True$
    \end{tabular}
  \end{center}
\end{enumerate}
\end{defn}

كما يسهل أن نرى، فإن أي~!!{formula}~$!A$ لا تحتوي إلا~$\lnot$ و
$\land$ و$\lor$ تأخذ قيمة الصدق~$\Undef$ إذا أسندت~$\Undef$ إلى
جميع~!!{propositional variable}s فيها. ولذلك، مثلًا، لا يكون
التحصيلان المنطقيان الكلاسيكيان~$p \lor \lnot p$ و
$\lnot(p \land \lnot p)$ تحصيلين منطقيين في~$\LogLuk[3]$، لأن
$\pValue{v}(!A) = \Undef$ كلما كان~$\pAssign v(p) = \Undef$.

في~!!{valuation}s التي يكون فيها~$\pAssign v(p) = \True$ أو~$\False$،
تتفق~$\pValue v(!A)$ مع قيمة صدقها الكلاسيكية.

\begin{prop}
  إذا كان~$\pAssign v(p) \in \{\True, \False\}$ لكل~$p$ في~$!A$،
  فإن~$\pValue v[\LogLuk[3]](!A) = \pValue v[\LogCL](!A)$.
\end{prop}

\begin{prob}\label{mvl:thr:luk:prob:luk-iff} افترض أننا نعرّف
  $\pValue v(!A \liff !B) = \pValue v((!A \lif !B) \land (!B \lif
  !A))$ في~$\LogLuk[3]$. فما جدول صدق~$\liff$؟
\end{prob}

كثير من التحصيلات المنطقية الكلاسيكية \emph{تحصيلات} منطقية أيضًا
في~\LogLuk[3]، مثل~$\lnot p \lif (p \lif q)$. وكما في المنطق
الكلاسيكي، يمكننا التحقق من ذلك بجداول الصدق:
\begin{center}
\begin{tabular}{cc|cccccc}
  $p$ & $q$ & $\lnot$ & $p$ & $\lif$ & $\smash{(}p$ & $\lif$ & $q\smash{)}$ \\ \hline
  \True & \True & \False & \True & \True & \True & \True & \True \\
  \True & \Undef & \False & \True & \True & \True & \Undef & \Undef \\
  \True & \False & \False & \True & \True & \True & \False & \False \\
  \Undef & \True & \Undef & \Undef & \True & \Undef & \True & \True \\
  \Undef & \Undef & \Undef & \Undef & \True & \Undef & \True & \Undef \\
  \Undef & \False & \Undef & \Undef & \True & \Undef & \Undef & \False \\
  \False & \True & \True & \False & \True & \False & \True & \True \\
  \False & \Undef & \True & \False & \True & \False & \True & \Undef \\
  \False & \False & \True & \False & \True & \False & \True & \False \\
\end{tabular}
\end{center}

\begin{prob}
  بين أن الآتي تحصيلات منطقية في~\LogLuk[3]:
  \begin{enumerate}
    \item $p \lif (q \lif p)$
    \item\label{mvl:thr:luk:prob:luk-taut-2} $\lnot(p \land q) \liff (\lnot p \lor \lnot q)$
    \item\label{mvl:thr:luk:prob:luk-taut-3} $\lnot(p \lor q) \liff (\lnot p \land \lnot q)$
  \end{enumerate}
  (في \olref[mvl][thr][luk]{prob:luk-taut-2} و
  \olref[mvl][thr][luk]{prob:luk-taut-3}، اعتبر~$!A \liff !B$
  اختصارًا لـ$(!A \lif !B) \land (!B \lif !A)$، أو ارجع إلى حلك
  لـ\olref[mvl][thr][luk]{prob:luk-iff}.)
\end{prob}

\begin{prob}
  بين أن التحصيلات المنطقية الكلاسيكية الآتية ليست تحصيلات منطقية
  في~\LogLuk[3]:
  \begin{enumerate}
    \item $(\lnot p \land p) \lif q$
    \item $((p \lif q) \lif p) \lif p$
    \item $(p \lif (p \lif q)) \lif (p \lif q)$
  \end{enumerate}
\end{prob}

قد يظن المرء لذلك أنه حتى إذا لم تكن جميع التحصيلات المنطقية
الكلاسيكية تحصيلات في~$\LogLuk[3]$، فلا بد أن تأخذ على الأقل إما
القيمة~$\True$ أو القيمة~$\Undef$ في كل~!!{valuation}. وهذا غير صحيح؛
إذ يعطي الصيغة الآتية مثالًا مضادًا:
\[
  \lnot(p \lif \lnot p) \lor \lnot(\lnot p \lif p)
\]
فهي تأخذ~$\False$ إذا أخذ~$p$ القيمة~$\Undef$.

\begin{prob}
  أي العلاقات الآتية تصح في منطق~\L ukasiewicz؟ أعط جدول صدق لكل منها.
  \begin{enumerate}
    \item $p, p \lif q \Entails q$
    \item $\lnot\lnot p \Entails p$
    \item $p \land q \Entails p$
    \item $p \Entails p \land p$
    \item $p \Entails p \lor q$
  \end{enumerate}
\end{prob}

أمل~\L ukasiewicz أن يبني منطقًا للإمكان على أساس نسقه الثلاثي القيم،
بإدخال رابط ذي موضع واحد~$\Diamond !A$ (بمعنى ``$!A$ ممكنة'') ورابط
موافق~$\Box !A$ (بمعنى ``$!A$ ضرورية''):
\begin{center}
  \begin{tabular}{c|c}
    $\tf{\Diamond}$ & \\
    \hline
    $\True$ & $\True$ \\
    $\Undef$ & $\True$ \\
    $\False$ & $\False$
  \end{tabular}
  \quad
  \begin{tabular}{c|c}
    $\tf{\Box}$ & \\
    \hline
    $\True$ & $\True$ \\
    $\Undef$ & $\False$ \\
    $\False$ & $\False$
  \end{tabular}
\end{center}
وبعبارة أخرى، يكون~$p$ ممكنًا إذا وفقط إذا لم يتعين كاذبًا بالفعل؛
ويكون~$p$ ضروريًا إذا وفقط إذا تعين صادقًا بالفعل.

\begin{prob}
  بين أن~$\Box p \liff \lnot\Diamond \lnot p$ و
  $\Diamond p \liff \lnot \Box \lnot p$ تحصيلان منطقيان في
  $\LogLuk[3]$ بعد توسيعه بجدولي صدق~$\Box$ و$\Diamond$.
\end{prob}

غير أن أوجه قصور هذا المنطق الجهي المقترح سرعان ما ظهرت: كيفما آلت
الأمور، لا يمكن أن يتبين أبدًا أن~$p \land \lnot p$ صادق. ولذلك، حتى
إذا لم يكن متعينًا الآن (فكان غير محدد)، ينبغي أن يعد مستحيلًا؛ أي
ينبغي أن تكون~$\lnot \Diamond(p \land \lnot p)$ تحصيلًا منطقيًا.
لكن إذا كان~$\pAssign v(p) = \Undef$، فإن
$\pValue v(\lnot \Diamond(p \land \lnot p)) = \Undef$. ومع أن
\L ukasiewicz كان محقًا في أن قيمتي صدق لا تكفيان لاستيعاب تمييزات
جهية كالإمكان والضرورة، فإن إدخال قيمة صدق ثالثة لا يكفي أيضًا.

\end{document}
